\documentclass[journal]{IEEEtran}
% Every number, table and figure in this file was generated from the per-seed logs by
% https://github.com/Abraheem13/tierfed and is written here as plain text.
\usepackage[utf8]{inputenc}
\usepackage[T1]{fontenc}
\usepackage{amsmath,amssymb,amsthm,booktabs,graphicx,xcolor,etoolbox,url,hyperref}
\usepackage{cite}
\hypersetup{hidelinks, pdftitle={Supplementary Material: Quota-Constrained Multi-Rate Synchronization in Federated Learning: Gain Bound, Peak-Load Control, and Schedule Leakage}, pdfauthor={Abraheem Rashid, Mehwish Mahmood, Pengpeng Hu}}
\newcommand{\tierfed}{TierFed}
\newcommand{\Gh}{\hat\Gamma}
\hyphenation{Fed-Avg Fed-Prox Fed-Per Fed-LAMA Ada-Comm Tier-Fed Res-Net}
\newsavebox\tabbox
\newcommand{\fit}[1]{\sbox\tabbox{#1}\ifdim\wd\tabbox>\linewidth\resizebox{\linewidth}{!}{\usebox\tabbox}\else\usebox\tabbox\fi}
\renewcommand{\thetable}{S-\Roman{table}}
\renewcommand{\thefigure}{S-\arabic{figure}}
\renewcommand{\theequation}{S\arabic{equation}}
\newcommand{\fig}[2]{\includegraphics[width=#2]{figures/#1.pdf}}
\newtheorem{lemma}{Lemma}
\renewcommand{\thesection}{S\arabic{section}}
\renewcommand{\thesubsection}{\thesection.\arabic{subsection}}
\AtBeginDocument{\def\UrlBigBreakPenalty{10000}}
\begin{document}
\bstctlcite{IEEEexample:BSTcontrol}
\title{Supplementary Material: Quota-Constrained Multi-Rate Synchronization in Federated Learning: Gain Bound, Peak-Load Control, and Schedule Leakage}
\author{Abraheem~Rashid, Mehwish~Mahmood, and Pengpeng~Hu
\thanks{\emph{(Corresponding author: Pengpeng Hu.)}}
\thanks{A.~Rashid is with the Department of Computer Science, Brunel University London, Uxbridge UB8~3PH, U.K. (e-mail: 2546393@brunel.ac.uk).}
\thanks{M.~Mahmood is with Queen's University Belfast, Belfast BT7~1NN, U.K.}
\thanks{P.~Hu is with the Department of Materials, The University of Manchester, Manchester M13~9PL, U.K. (e-mail: pengpeng.hu@manchester.ac.uk).}}
\maketitle

This supplement contains the proofs of the two propositions and of the phase condition (Section~\ref{supp:proofs}), details of the privacy accounting (Section~\ref{supp:dp}), the sensitivity of the primary contrasts to the rule that selects the reported round (Section~\ref{supp:selection}), a paired comparison on the main benchmark (Section~\ref{supp:paired}), the privacy planner on a grid of its parameters (Section~\ref{supp:planner}), details of the two synchronization steps (Section~\ref{supp:sync}) and of the clinical-center benchmark (Section~\ref{supp:centers}), the experiment families of the main protocol (Section~\ref{supp:families}), additional experiments that were read at the final round (Section~\ref{supp:legacy}), and the location of the released data (Section~\ref{supp:data}). Equation, table, and section numbers without the prefix S refer to the article. Abbreviations are those of the article; AUPRC denotes the area under the precision--recall curve.

\section{Proofs}
\label{supp:proofs}
The notation and the assumptions are those of Section~\mbox{IV-B} of the article: local SGD with step size $\eta$, $s$ local steps per epoch at every client, equal weights $1/N$, block-separable client objectives $F_k(\theta)=\sum_gF_{k,g}(\theta_g)$ with $L$-smooth blocks, unbiased stochastic gradients with variance at most $\sigma_g^2$ on tier $g$ that are independent across clients and steps, and heterogeneity bounded by $\Gamma_g$. Because the blocks are separable, each tier evolves independently of the others, and we drop the index $g$ where no confusion can arise.

\subsection{Proposition 1(a): Drift of the Client Copies}
\begin{lemma}[Drift from a common value]
\label{lem:drift}
Let all clients start from a common value $\theta^0$ of a tier and take $t$ steps of SGD on their own objectives, $\theta_k^{\tau+1}=\theta_k^\tau-\eta(\nabla F_k(\theta_k^\tau)+\zeta_k^\tau)$, where $\zeta_k^\tau$ has zero mean given the past and $\mathbb E\|\zeta_k^\tau\|^2\le\sigma^2$. Let $\bar\theta^\tau=\frac1N\sum_k\theta_k^\tau$ and $V_\tau=\frac1N\sum_k\mathbb E\|\theta_k^\tau-\bar\theta^\tau\|^2$. If $\eta\le1/(\sqrt6\,Lt)$, then
\begin{equation}
V_t\le6\eta^2\big(t\sigma^2+t^2\Gamma\big).
\label{eq:s-drift}
\end{equation}
\end{lemma}
\begin{proof}
Write $a_k^u=\nabla F_k(\theta_k^u)-\nabla F_k(\bar\theta^u)$ and $h_k^u=\nabla F_k(\bar\theta^u)-\nabla F(\bar\theta^u)$, and denote averages over the clients by a bar. Since $F=\frac1N\sum_kF_k$, the difference between the update of client $k$ and the mean update at step $u$ is $(\zeta_k^u-\bar\zeta^u)+(a_k^u-\bar a^u)+h_k^u$, so that
\[
\theta_k^\tau-\bar\theta^\tau=-\eta\sum_{u<\tau}\big[(\zeta_k^u-\bar\zeta^u)+(a_k^u-\bar a^u)+h_k^u\big].
\]
Since $\|x+y+z\|^2\le3(\|x\|^2+\|y\|^2+\|z\|^2)$, we may bound the three sums separately. The noise terms are martingale differences, so their cross terms vanish, and $\frac1N\sum_k\|x_k-\bar x\|^2\le\frac1N\sum_k\|x_k\|^2$ gives $\frac1N\sum_k\mathbb E\|\sum_{u<\tau}(\zeta_k^u-\bar\zeta^u)\|^2\le\tau\sigma^2$. For the other two sums, we use $\|\sum_{u<\tau}x_u\|^2\le\tau\sum_{u<\tau}\|x_u\|^2$; the smoothness of $F_k$ gives $\frac1N\sum_k\|a_k^u-\bar a^u\|^2\le L^2\frac1N\sum_k\|\theta_k^u-\bar\theta^u\|^2$, and the heterogeneity bound gives $\frac1N\sum_k\|h_k^u\|^2\le\Gamma$. Hence
\[
V_\tau\le3\eta^2\big(\tau\sigma^2+\tau^2\Gamma\big)+3\eta^2L^2\tau\sum_{u<\tau}V_u .
\]
Let $M_\tau=6\eta^2(\tau\sigma^2+\tau^2\Gamma)$, which is nondecreasing in $\tau$. Then $V_0=0\le M_0$, and if $V_u\le M_u$ for all $u<\tau\le t$, then $V_\tau\le\frac12M_\tau+3\eta^2L^2\tau^2M_\tau\le M_\tau$, because $3\eta^2L^2t^2\le\frac12$. The claim follows by induction.
\end{proof}
Part (a) of Proposition~1 is Lemma~\ref{lem:drift} with $t=ms$: after $m$ local epochs from a common value, $V\le6\eta^2(ms\,\sigma_g^2+m^2s^2\Gamma_g)\le6\eta^2s^2m^2(\Gamma_g+\sigma_g^2)$, and the condition on the step size holds for all $m\le K_{\max}$ when $\eta\le1/(\sqrt6LK_{\max}s)$, so (11) holds with $c_2=6$ for any $c_0\le1/\sqrt6$. Under the restart step, the clients' copies of tier $g$ leave a common value at two points. The first is the start of a round in which the tier is due, by (1); one local epoch then follows before (3) makes them equal again ($m=1$). The second is that adoption, after which they evolve for the $K_g-1$ rounds in which the tier is not due ($m\le K_g-1$). Therefore, $m\le\max(1,K_g-1)$.

\subsection{Proposition 1(b): The Server Copy}
In a round $r_j$ in which tier $g$ is due, every client starts from the server copy $\theta_g^{(r_j)}$ by (1) and takes $s$ steps of SGD on $F_{k,g}$, and the server averages the results by (2). This is one round of federated averaging with $s$ local steps on $F_g=\frac1N\sum_kF_{k,g}$. Between $r_j$ and $r_{j+1}$, the server copy does not change. What the clients do with their own copies of the tier in those rounds has no effect on round $r_{j+1}$: the copies are overwritten by (1), SGD carries no state from one step to the next, and, by block separability, the gradient on tier $g$ does not depend on the other tiers. The server copies $\theta_g^{(r_1)},\theta_g^{(r_2)},\dots$ are therefore exactly the iterates of federated averaging with $s$ local steps, full participation, and equal weights, run for $J_g$ rounds. For this algorithm, with $L$-smooth client objectives, unbiased stochastic gradients of variance at most $\sigma_g^2$, and gradient dissimilarity bounded by $\Gamma_g$, the standard nonconvex analysis (for instance, that of~\cite{wang2020fednova} with equal numbers of local steps, in which case its surrogate objective coincides with $F_g$; see also~\cite{karimireddy2020,wang2021coop}) yields, for $\eta\le c_0/(Ls)$ with a sufficiently small absolute constant $c_0\le1/\sqrt6$ and absolute constants $c_1$ and $c_2\ge6$,
\begin{align*}
\min_{j\le J_g}\mathbb E\big\|\nabla F_g(\theta_g^{(r_j)})\big\|^2&\le\frac{c_1D_g}{\eta sJ_g}+c_1\eta L\frac{\sigma_g^2}{N}\\
&\quad+c_2\eta^2L^2s^2\big(\Gamma_g+\sigma_g^2\big),
\end{align*}
which is (12), since $\nabla_gF(\theta)=\nabla F_g(\theta_g)$. The three terms have the usual origins. The first is the progress of $sJ_g$ gradient steps from an initial gap $D_g$. The second is the variance of the averaged stochastic gradient: the average $\frac1N\sum_k\zeta_k$ of $N$ independent noise vectors with $\mathbb E\|\zeta_k\|^2\le\sigma_g^2$ has $\mathbb E\|\frac1N\sum_k\zeta_k\|^2\le\sigma_g^2/N$, and, for a weighted average $\sum_kp_k\zeta_k$, the bound is $\sigma_g^2\sum_kp_k^2$, which is at least $\sigma_g^2/N$ and equals it only for equal weights. The third is the drift of the local iterates within one round, which is bounded as in Lemma~\ref{lem:drift} with $t\le s$. We make no claim about the values of the constants in (12).

\emph{Carry-over step.} If (1) is skipped after the first round, the clients' copies of tier $g$ evolve from the last aggregate for $K_gs$ steps before they are averaged again. This is local SGD with a period of $K_gs$ steps~\cite{stich2019,wang2021coop}. Its bound has three terms: the progress term $c_1D_g/(\eta sR)$, because all $sR$ local steps now contribute to the server copy; the same variance term; and the drift term $c_2\eta^2L^2s^2K_g^2(\Gamma_g+\sigma_g^2)$, which follows from Lemma~\ref{lem:drift} with $t\le K_gs$ and requires $\eta\le c_0/(LK_{\max}s)$.

\emph{What the scheduler measures.} The numerator of $\Gh_g$ in (10) is the dispersion of the uploaded slices of tier $g$. Under the restart step, every uploaded slice is $s$ steps away from a common value, so, by (\ref{eq:s-drift}), the expectation of this dispersion is at most $6\eta^2(s\sigma_g^2+s^2\Gamma_g)$ whatever the interval. The noise and the heterogeneity enter with different powers of $s$, and the denominator is the squared norm of the tier, so $\Gh_g$ is a proxy for $\Gamma_g+\sigma_g^2$ and not an estimate of it. Under the carry-over step, the slices are $K_gs$ steps away from a common value, and the same bound grows by a factor between $K_g$ (noise) and $K_g^2$ (heterogeneity). Section~\ref{supp:sync} shows both behaviors in the logs.

\subsection{Proposition 2}
\label{supp:cuberoot}
Write $x_g=K_g>0$. The objective $f(x)=\sum_g\Gamma_gx_g^2$ is strictly convex, the constraint $h(x)=\sum_gw_g/x_g-\rho\le0$ is convex on $x>0$, and Slater's condition holds, so the Karush--Kuhn--Tucker conditions are necessary and sufficient and the minimizer is unique. Since $f$ increases and $h$ decreases in every coordinate, the constraint is active: if $h(x)<0$, some $x_g$ can be lowered, which lowers $f$. Stationarity, $2\Gamma_gx_g=\lambda w_g/x_g^2$, gives $x_g=\big(\lambda w_g/(2\Gamma_g)\big)^{1/3}$. Substitution into $\sum_gw_g/x_g=\rho$ gives $(\lambda/2)^{1/3}=S/\rho$ with $S=\sum_hw_h^{2/3}\Gamma_h^{1/3}$, hence $K_g^\star=(w_g/\Gamma_g)^{1/3}S/\rho$ and $f(K^\star)=\frac{S^2}{\rho^2}\sum_g\Gamma_g(w_g/\Gamma_g)^{2/3}=S^3/\rho^2$. H\"older's inequality with exponents $3/2$ and $3$, applied to $a_g=w_g^{2/3}$ and $b_g=\Gamma_g^{1/3}$, gives $S\le(\sum_gw_g)^{2/3}(\sum_g\Gamma_g)^{1/3}=(\sum_g\Gamma_g)^{1/3}$, with equality if and only if $a_g^{3/2}\propto b_g^{3}$, that is, $w_g\propto\Gamma_g$. The uniform point $x_g=1/\rho$ satisfies the constraint with equality and costs $\sum_g\Gamma_g/\rho^2$. Finally, the integer feasible set is contained in the real one, so the integer optimum satisfies $J_{\mathrm{int}}\ge J^\star$; when $1/\rho\in\{1,\dots,K_{\max}\}$, the uniform period is integer-feasible, and $1-J_{\mathrm{int}}/J_{\mathrm u}\le1-J^\star/J_{\mathrm u}=\Psi$. The quantity $\Psi$ is invariant under $\Gamma\mapsto c\Gamma$, because $S^3$ and $\sum_g\Gamma_g$ are both linear in $c$.\hfill\IEEEQED

\subsection{Staggered Phases}
After warm-up, tier $a$ is due in the rounds $r=T_w+\phi_a+nK_a>T_w$ and tier $b$ in the rounds $r=T_w+\phi_b+mK_b>T_w$, for integers $n,m\ge0$ ($n\ge1$ when $\phi_a=0$ and $m\ge1$ when $\phi_b=0$). A common round requires $nK_a-mK_b=\phi_b-\phi_a$. This linear Diophantine equation is solvable if and only if $\gcd(K_a,K_b)$ divides $\phi_b-\phi_a$, and, when it is solvable, it has infinitely many solutions with $n,m\ge1$. Over an unbounded horizon, the two sets of due rounds are therefore disjoint if and only if $\phi_a-\phi_b\not\equiv0\pmod{\gcd(K_a,K_b)}$; the first common round, if there is one, falls within $\mathrm{lcm}(K_a,K_b)$ rounds after warm-up.

\section{Differential-Privacy Accounting}
\label{supp:dp}
Clipping each client's update to $C$ makes the $\ell_2$ sensitivity of the sum to the addition or removal of one client equal to $C$, and the fixed denominator $qN$ of (17) does not depend on the realized cohort. With Poisson sampling at rate $q$, each release is a subsampled Gaussian mechanism, whose R\'enyi differential privacy (DP) at integer order $\alpha$ is $\frac1{\alpha-1}\log\sum_{k=0}^{\alpha}\binom{\alpha}{k}(1-q)^{\alpha-k}q^k\exp\frac{k^2-k}{2z^2}$~\cite{mironov2019sgm}; for $q=1$, it is $\alpha/(2z^2)$~\cite{mironov2017}. Composition over the $R_{\mathrm{tx}}$ planned releases is additive, and the conversion to $(\varepsilon,\delta)$-DP~\cite{mironov2017} is minimized over $\alpha\in\{2,\dots,256\}$ for $q<1$ and over $\alpha\in\{1.1,1.2,\dots,10.9\}\cup\{12,13,\dots,255\}$ for $q=1$. We evaluate the sum in log space. The implementation agrees with the TensorFlow Privacy MNIST reference~\cite{tfprivacy} ($\varepsilon\approx3.0$ at $q=256/60000$, $z=1.1$, 60 epochs) to within 0.05. The accountant is valid only if the set of release rounds is independent of the private data, which is why we run only FedPer, the fixed periods, one-shot federation, and the planner under DP in the article. Model selection on held-out data is not covered by the budget. The DP tables of Section~\ref{supp:legacy} therefore report final-round values, and Table~8 of the article gives the final round as its primary value and the validation-selected value as an uncharged reference only.

\emph{Exact accountant at full participation.} At $q=1$, the $R_{\mathrm{tx}}$ releases are Gaussian mechanisms without subsampling, whose exact privacy is $\mu$-Gaussian DP with $\mu=\sqrt{R_{\mathrm{tx}}}/z$~\cite{dong2022gdp}. The corresponding $(\varepsilon,\delta)$ curve~\cite{balle2018} is
\begin{equation}
\delta(\varepsilon)=\Phi\Big(-\frac{\varepsilon}{\mu}+\frac{\mu}{2}\Big)-e^{\varepsilon}\,\Phi\Big(-\frac{\varepsilon}{\mu}-\frac{\mu}{2}\Big),
\label{eq:gdp}
\end{equation}
where $\Phi$ is the standard normal distribution function. Like (18), it depends on $R_{\mathrm{tx}}$ only through $R_{\mathrm{tx}}/z^2$. Calibrated on (\ref{eq:gdp}) instead of (18), the noise multiplier of every method would be lower by the same factor: 24\%, 20\%, 17\%, and 13\% for $\varepsilon^\star=1$, 2, 4, and 8, respectively. The article uses the R\'enyi accountant, so every $\varepsilon$ that it reports is an upper bound.

\section{Sensitivity to the Round-Selection Rule}
\label{supp:selection}
The primary rule of the article selects the first round, $\hat r$, that maximizes the mean validation AUROC, $A^{\mathcal V}$, of (7). Under a sparse schedule, the test AUROC rises at each synchronization and falls afterward, so this rule can select the peak that follows a synchronization. On a flat validation curve, a small numerical difference can move the selected round (Section~\mbox{VI-A}). Table~\ref{tab:selection} repeats both primary contrasts under two further rules that are computed from the same per-round logs: the first round that maximizes the trailing three-round mean of the validation AUROC, and the mean test AUROC over the last 10 rounds, which involves no selection. The first rule of the table is the primary rule of the article. Its differences are those of Table~\ref{tab:paired} and of Table~5 of the article; every bootstrap reads the seeds in ascending order, so the intervals agree between tables.

Under none of the three rules, and in no setting, is \tierfed{} ahead of the best uniform period by at least the relevance threshold of 0.005 with an interval above zero. The comparison with FedPer depends on the rule. Under the smoothed rule, \tierfed{} is below FedPer by more than the threshold, with an interval below zero, in two settings. Over the last 10 rounds, it is above FedPer, with an interval above zero, in five of the seven settings, because full-rate training has overfitted by then (Section~\mbox{VI-B}).

\begin{table*}[t]
\caption{The two primary contrasts under three rules for the reported value, computed from the same runs: mean paired difference in AUROC with its bootstrap 95\% interval. The best uniform period is the one chosen in Table~5 of the article.}
\label{tab:selection}
\centering\small
\fit{\begin{tabular}{@{}llcc@{}}
\toprule
Setting, seeds, best uniform period & Reported value & \tierfed{} $-$ FedPer [95\% CI] & \tierfed{} $-$ best uniform [95\% CI]\\
\midrule
Specialty groups, 10 seeds, $K=4$ & First maximum of $A^{\mathcal V}$ (article) & \ensuremath{-}0.0007 [\ensuremath{-}0.0033, \ensuremath{+}0.0018] & \ensuremath{+}0.0000 [\ensuremath{+}0.0000, \ensuremath{+}0.0000]\\
 & Maximum of the three-round mean of $A^{\mathcal V}$ & \ensuremath{-}0.0030 [\ensuremath{-}0.0056, \ensuremath{-}0.0005] & \ensuremath{+}0.0000 [\ensuremath{+}0.0000, \ensuremath{+}0.0000]\\
 & Mean of the last 10 rounds & \ensuremath{+}0.0085 [\ensuremath{+}0.0051, \ensuremath{+}0.0118] & \ensuremath{+}0.0000 [\ensuremath{+}0.0000, \ensuremath{+}0.0000]\\
\midrule
Dirichlet $\alpha=0.1$, 5 seeds, $K=8$ & First maximum of $A^{\mathcal V}$ (article) & \ensuremath{-}0.0010 [\ensuremath{-}0.0182, \ensuremath{+}0.0215] & \ensuremath{-}0.0113 [\ensuremath{-}0.0270, \ensuremath{+}0.0040]\\
 & Maximum of the three-round mean of $A^{\mathcal V}$ & \ensuremath{-}0.0091 [\ensuremath{-}0.0224, \ensuremath{+}0.0058] & \ensuremath{-}0.0099 [\ensuremath{-}0.0289, \ensuremath{+}0.0004]\\
 & Mean of the last 10 rounds & \ensuremath{-}0.0052 [\ensuremath{-}0.0208, \ensuremath{+}0.0084] & \ensuremath{-}0.0096 [\ensuremath{-}0.0219, \ensuremath{-}0.0008]\\
\midrule
Dirichlet $\alpha=0.3$, 5 seeds, $K=4$ & First maximum of $A^{\mathcal V}$ (article) & \ensuremath{+}0.0354 [\ensuremath{+}0.0075, \ensuremath{+}0.0788] & \ensuremath{+}0.0000 [\ensuremath{+}0.0000, \ensuremath{+}0.0000]\\
 & Maximum of the three-round mean of $A^{\mathcal V}$ & \ensuremath{+}0.0342 [\ensuremath{+}0.0089, \ensuremath{+}0.0726] & \ensuremath{+}0.0000 [\ensuremath{+}0.0000, \ensuremath{+}0.0000]\\
 & Mean of the last 10 rounds & \ensuremath{+}0.0100 [\ensuremath{+}0.0002, \ensuremath{+}0.0221] & \ensuremath{+}0.0000 [\ensuremath{+}0.0000, \ensuremath{+}0.0000]\\
\midrule
Dirichlet $\alpha=1$, 5 seeds, $K=4$ & First maximum of $A^{\mathcal V}$ (article) & \ensuremath{+}0.0029 [\ensuremath{-}0.0049, \ensuremath{+}0.0112] & \ensuremath{+}0.0000 [\ensuremath{+}0.0000, \ensuremath{+}0.0000]\\
 & Maximum of the three-round mean of $A^{\mathcal V}$ & \ensuremath{+}0.0038 [\ensuremath{-}0.0086, \ensuremath{+}0.0161] & \ensuremath{+}0.0000 [\ensuremath{+}0.0000, \ensuremath{+}0.0000]\\
 & Mean of the last 10 rounds & \ensuremath{+}0.0140 [\ensuremath{+}0.0041, \ensuremath{+}0.0239] & \ensuremath{+}0.0000 [\ensuremath{+}0.0000, \ensuremath{+}0.0000]\\
\midrule
Dirichlet $\alpha=10$, 5 seeds, $K=4$ & First maximum of $A^{\mathcal V}$ (article) & \ensuremath{+}0.0000 [\ensuremath{-}0.0021, \ensuremath{+}0.0028] & \ensuremath{+}0.0000 [\ensuremath{+}0.0000, \ensuremath{+}0.0000]\\
 & Maximum of the three-round mean of $A^{\mathcal V}$ & \ensuremath{-}0.0125 [\ensuremath{-}0.0168, \ensuremath{-}0.0090] & \ensuremath{+}0.0000 [\ensuremath{+}0.0000, \ensuremath{+}0.0000]\\
 & Mean of the last 10 rounds & \ensuremath{+}0.0017 [\ensuremath{-}0.0005, \ensuremath{+}0.0043] & \ensuremath{+}0.0000 [\ensuremath{+}0.0000, \ensuremath{+}0.0000]\\
\midrule
ResNet backbone, 10 seeds, $K=4$ & First maximum of $A^{\mathcal V}$ (article) & \ensuremath{-}0.0029 [\ensuremath{-}0.0057, \ensuremath{-}0.0002] & \ensuremath{-}0.0023 [\ensuremath{-}0.0061, \ensuremath{+}0.0011]\\
 & Maximum of the three-round mean of $A^{\mathcal V}$ & \ensuremath{-}0.0055 [\ensuremath{-}0.0088, \ensuremath{-}0.0025] & \ensuremath{+}0.0038 [\ensuremath{-}0.0004, \ensuremath{+}0.0074]\\
 & Mean of the last 10 rounds & \ensuremath{+}0.0078 [\ensuremath{+}0.0047, \ensuremath{+}0.0113] & \ensuremath{-}0.0062 [\ensuremath{-}0.0077, \ensuremath{-}0.0047]\\
\midrule
Clinical centers, 20 seeds, $K=4$ & First maximum of $A^{\mathcal V}$ (article) & \ensuremath{-}0.0085 [\ensuremath{-}0.0161, \ensuremath{-}0.0010] & \ensuremath{+}0.0000 [\ensuremath{+}0.0000, \ensuremath{+}0.0000]\\
 & Maximum of the three-round mean of $A^{\mathcal V}$ & \ensuremath{-}0.0059 [\ensuremath{-}0.0125, \ensuremath{+}0.0000] & \ensuremath{+}0.0000 [\ensuremath{+}0.0000, \ensuremath{+}0.0000]\\
 & Mean of the last 10 rounds & \ensuremath{+}0.0063 [\ensuremath{+}0.0012, \ensuremath{+}0.0114] & \ensuremath{+}0.0000 [\ensuremath{+}0.0000, \ensuremath{+}0.0000]\\
\bottomrule
\end{tabular}}
\end{table*}

\section{Paired Comparison on the Specialty Groups}
\label{supp:paired}
Table~\ref{tab:paired} compares \tierfed{} with every baseline of Table~3 of the article on the same 10 seeds. We omit the fixed period $(4,4)$, which is identical to \tierfed{}. The $p$-values are corrected by Holm's method over the listed contrasts.

\begin{table}[t]
\caption{Paired comparison of \tierfed{} with each baseline on the specialty groups: difference in AUROC at $\hat r$ with its bootstrap 95\% interval, Holm-corrected Wilcoxon $p$-value, number of seeds on which \tierfed{} is strictly higher, and difference in upload in MB up to $\hat r$ and over 60 rounds.}
\label{tab:paired}
\centering\small
\setlength{\tabcolsep}{3pt}\fit{\begin{tabular}{@{}lccccc@{}}
\toprule
\tierfed{} against & $\Delta$AUROC [95\% CI] & $p_{\mathrm{Holm}}$ & Higher & \shortstack{$\Delta$MB\\to $\hat r$} & \shortstack{$\Delta$MB,\\60 rounds}\\
\midrule
Local only & \ensuremath{+}0.0232 [\ensuremath{+}0.0187, \ensuremath{+}0.0271] & 0.020 & 10/10 & \ensuremath{+}45.7 & \ensuremath{+}65.3\\
FedAvg & \ensuremath{-}0.0014 [\ensuremath{-}0.0042, \ensuremath{+}0.0012] & 1.000 & 4/10 & \ensuremath{-}4.9 & \ensuremath{-}153.1\\
FedProx & \ensuremath{-}0.0014 [\ensuremath{-}0.0046, \ensuremath{+}0.0023] & 1.000 & 2/10 & \ensuremath{-}32.2 & \ensuremath{-}153.1\\
FedPer & \ensuremath{-}0.0007 [\ensuremath{-}0.0033, \ensuremath{+}0.0018] & 1.000 & 4/10 & \ensuremath{-}17.4 & \ensuremath{-}152.5\\
FedLAMA-style & \ensuremath{-}0.0019 [\ensuremath{-}0.0045, \ensuremath{+}0.0002] & 1.000 & 3/10 & \ensuremath{-}3.7 & \ensuremath{-}88.7\\
AdaComm & \ensuremath{+}0.0048 [\ensuremath{+}0.0021, \ensuremath{+}0.0077] & 0.088 & 9/10 & \ensuremath{+}18.1 & \ensuremath{+}36.2\\
LAG, equal weights & \ensuremath{+}0.0005 [\ensuremath{-}0.0028, \ensuremath{+}0.0042] & 1.000 & 4/10 & \ensuremath{-}71.1 & \ensuremath{-}153.1\\
LAG & \ensuremath{-}0.0018 [\ensuremath{-}0.0044, \ensuremath{+}0.0004] & 1.000 & 4/10 & \ensuremath{+}1.3 & \ensuremath{-}153.1\\
Fixed (8,8) & \ensuremath{+}0.0005 [\ensuremath{-}0.0015, \ensuremath{+}0.0023] & 1.000 & 7/10 & \ensuremath{+}9.8 & \ensuremath{+}25.4\\
Staggered (4,4) & \ensuremath{+}0.0024 [\ensuremath{+}0.0000, \ensuremath{+}0.0052] & 0.844 & 7/10 & \ensuremath{+}5.5 & \ensuremath{-}1.6\\
\bottomrule
\end{tabular}}
\end{table}

\section{The Privacy Planner on a Grid}
\label{supp:planner}
Table~\ref{tab:planner} gives the intervals that the planner of Section~\mbox{IV-D} selects by minimizing (19), computed with the implementation that the experiments use, for the tier shares of the MLP on the main benchmark, $R=60$, $T_w=5$, $K_{\max}=8$, $q=1$, and $\delta=10^{-5}$. The grid has seven values of the noise weight $\lambda$, five quotas, and four budgets $\varepsilon^\star\in\{1,2,4,8\}$. Because the ratio of noise energies in (19) does not depend on $\varepsilon^\star$ at $q=1$, the selected pair is the same for all four budgets, and the table therefore has one row per quota. The quota is applied to the realized payload fraction $\Pi_R$, warm-up included; quotas below $(5+6)/60\approx0.18$ are infeasible, since the sparsest pair $(8,8)$ already transmits in 11 of the 60 rounds. Without a binding quota ($\rho=1$), the planner returns equal intervals for every $\lambda$, and a larger $\lambda$, which weights the noise more, gives longer intervals. The only unequal pair is $(5,6)$ at $\rho=0.25$: there, $(4,4)$ realizes a payload fraction of 0.30 and is infeasible, and, among the feasible pairs, the planner gives the longer interval to the medium tier, which carries the larger share of the payload.

\begin{table}[t]
\caption{Intervals (slow, medium) selected by the privacy planner for each quota $\rho$ on the realized payload fraction and each noise weight $\lambda$; the selection is the same for $\varepsilon^\star\in\{1,2,4,8\}$.}
\label{tab:planner}
\centering\small
\fit{\begin{tabular}{@{}cccccccc@{}}
\toprule
Quota $\rho$ & $\lambda=0.1$ & $\lambda=0.25$ & $\lambda=0.5$ & $\lambda=1$ & $\lambda=2$ & $\lambda=4$ & $\lambda=10$\\
\midrule
1 & (2,2) & (2,2) & (2,2) & (3,3) & (4,4) & (4,4) & (5,5)\\
0.4 & (3,3) & (3,3) & (3,3) & (3,3) & (4,4) & (4,4) & (5,5)\\
0.3 & (4,4) & (4,4) & (4,4) & (4,4) & (4,4) & (4,4) & (5,5)\\
0.25 & (5,6) & (5,6) & (5,6) & (5,6) & (6,6) & (6,6) & (6,6)\\
0.2 & (7,7) & (7,7) & (7,7) & (7,7) & (7,7) & (7,7) & (7,7)\\
\bottomrule
\end{tabular}}
\end{table}

\section{The Two Synchronization Steps}
\label{supp:sync}
Section~\mbox{III-A} of the article defines the restart step and the carry-over step. The carry-over step is examined in Section~\mbox{VI-E}, and every other experiment of the main protocol uses the restart step. The two differ in one line of the round: under the carry-over step, a client does not replace its copy of a due tier with the server copy before it trains, except in the first round, so the tier that it uploads is the result of all local epochs since the tier's last synchronization. Full-rate methods are identical under both steps, because a tier that is due in every round is adopted from the server at the end of each round by (3); we ran FedPer under the carry-over step for one seed and reproduced its restart run exactly. The carry-over runs use full participation and no DP.

Table~\ref{tab:dispersion} reports the dispersion $\Gh_g$ of the uploaded slices, averaged over the synchronization rounds after warm-up, for uniform periods on the main benchmark. Under the restart step, it varies by a factor of at most 1.29 between $K=1$ and $K=8$ and shows no trend in $K$. Under the carry-over step, it is larger by factors of 7.5 and 5.3 at $K=4$ and of 26.9 and 11.9 at $K=8$ for the slow and the medium tier, respectively. This growth lies between the linear and the quadratic growth that Lemma~\ref{lem:drift} allows. A scheduler that reads $\Gh_g$ under the carry-over step therefore measures the interval that it chose earlier together with the disagreement of the clients. The variant of \tierfed{} with a rescaled dispersion, listed in Table~6 of the article, divides $\Gh_g$ by the square of the number of rounds since the tier's last synchronization to remove this dependence. On the specialty groups and at the clinical centers, both variants select $(4,4)$ in every seed under the carry-over step, and their runs are identical to those of the fixed period $(4,4)$ under that step; the ResNet is discussed in Section~\mbox{VI-E} of the article.

\begin{table}[t]
\caption{Dispersion $\Gh_g$ ($\times10^{-3}$) of the uploaded slices of the two shared tiers, averaged over the synchronization rounds after warm-up, for uniform periods $K$ on the specialty groups (MLP) under the two synchronization steps; $K=1$ is FedPer, for which the steps coincide; a dash marks a combination that was not run.}
\label{tab:dispersion}
\centering\small
\fit{\begin{tabular}{@{}ccccccc@{}}
\toprule
 & \multicolumn{3}{c}{Restart step} & \multicolumn{3}{c}{Carry-over step}\\
\cmidrule(lr){2-4}\cmidrule(lr){5-7}
Period $K$ & Seeds & Slow tier & Medium tier & Seeds & Slow tier & Medium tier\\
\midrule
1 & 10 & 4.81 & 7.40 & -- & -- & --\\
2 & 5 & 5.93 & 7.46 & -- & -- & --\\
4 & 10 & 5.70 & 6.93 & 10 & 42.90 & 36.67\\
6 & 5 & 5.10 & 7.22 & -- & -- & --\\
8 & 10 & 4.59 & 7.71 & 10 & 123.25 & 91.74\\
\bottomrule
\end{tabular}}
\end{table}

\section{The Clinical-Center Benchmark}
\label{supp:centers}
Fed-Heart-Disease~\cite{duterrail2022flamby} consists of the four processed collections of the UCI Heart Disease data~\cite{detrano1989}: Cleveland (303 patients), Hungary (294), Switzerland (123), and Long Beach (200), for 920 patients in total. The FLamby release removes missing values and retains 740 of these patients; we keep all 920. The label is the presence of angiographic disease (a diagnosis value above zero). The 13 clinical attributes are encoded as 36 features. The six numeric attributes are standardized; a missing value is replaced by the median and marked by an indicator feature, because some attributes are mostly missing outside the Cleveland collection. The seven categorical attributes are one-hot encoded, with a missing value as a category of its own. As on the main benchmark, the medians and the scaling use statistics of the pooled data and do not involve the labels. Each collection is one client. A fifth of the patients are set aside and enter no result, as on the main benchmark, and the remaining 736 patients are split per center into training, validation, and test sets (70/10/20\%). The MLP of the main benchmark has 29,634 parameters on this input, of which the slow, medium, and fast tiers hold 16.0\%, 83.6\%, and 0.4\%, respectively. All hyperparameters are those of Table~2 of the article, including the batch size of 32 that is used at the centers. The methods are those of the main comparison, except that we omit LAG with equal weights and add the uniform period $K=2$; the statistical procedure is the same, here with 20 seeds. We add $K=2$ so that the matched-bytes contrast has four candidate periods ($K\in\{1,2,4,8\}$).

Because the centers are small, the test split of a center holds between about 19 and 50 patients, and, in 30\% of the runs, one center has a test split with a single class and is removed from the average of that run for every method alike. The results of Table~4 of the article should therefore be read as a check that the behavior observed on the specialty groups carries over to clients that are real institutions, and not as a clinical evaluation.

\section{Experiment Families of the Main Protocol}
\label{supp:families}
Table~\ref{tab:families} lists the 946 runs of the main protocol of the article by family, with the numbers of configurations, seeds, and runs counted from the logs. A configuration is a method with its parameters in one setting; the timing families count the runs of both worlds. The single-seed FedPer check of Section~\ref{supp:sync} is not counted.

\begin{table*}[t]
\caption{Experiment families of the main protocol: data and clients, number of configurations, seeds per configuration, and runs.}
\label{tab:families}
\centering\small
\fit{\begin{tabular}{@{}llccc@{}}
\toprule
Family & Data and clients & Configurations & Seeds & Runs\\
\midrule
Main comparison & 20 specialty groups & 13 & 10 & 130\\
Reference bounds & pooled data, per-group test splits & 2 & 10 & 20\\
Accuracy--upload frontier & 20 specialty groups & 8 & 5 & 40\\
Peak-aware variants & 20 specialty groups & 3 & 5 & 15\\
Label-skew sweep & Dirichlet $\alpha\in\{0.1,0.3,1,10\}$ & 20 & 5 & 100\\
ResNet backbone & 20 specialty groups & 5 & 10 & 50\\
Clinical centers & Fed-Heart-Disease, 4 centers & 13 & 20 & 260\\
Carry-over synchronization step & specialty groups, centers, ResNet & 13 & 5--20 & 165\\
Matched-$\varepsilon$ DP ($q=1$) & 20 specialty groups & 5 & 10 & 50\\
Timing side channel, MLP & clean or shifted target group & 5 & 10 per world & 100\\
Timing side channel, ResNet & clean or shifted target group & 1 & 8 per world & 16\\
\midrule
Total &  &  &  & 946\\
\bottomrule
\end{tabular}}
\end{table*}

\section{Additional Experiments Read at the Final Round}
\label{supp:legacy}
The experiments of this section cover settings that the main protocol of the article does not: sweeps of the baselines' parameters, ablations, further backbones, imaging data, scaling, a transport check, and sweeps of the privacy budget. We ran them with a different configuration from that of the article, and their values are not comparable with those of its tables. Each client's patients are split 80/20\% into training and test sets without a validation set, and the AUROC is read at the final round, which can favor sparse schedules over full-rate methods (Section~\mbox{VI-B}). The fixed schedule of these tables synchronizes the slow tier every five rounds and the medium tier every two rounds from round 1, without warm-up, and LAG averages with equal client weights. The greedy solver of Table~\ref{tab:periods} repeatedly doubles the interval of the tier for which the doubling raises the objective of (13) least per unit of payload saved, until the quota is met, so its intervals are powers of two. All tables of this section are generated from the per-seed logs by the script \texttt{src/analysis/collect.py} of the released repository.

Table~\ref{tab:periods} reports the sweep of the fixed period and the comparison of the two solvers; Table~\ref{tab:baselines-legacy}, the sweeps of the baselines and the methods with an SGD local solver; Table~\ref{tab:lambda}, the sensitivity of the privacy planner to its noise weight; Table~\ref{tab:exposure}, the exposure of the transmitted updates to an update-inversion attack~\cite{geiping2020} and to a loss-threshold membership-inference attack~\cite{yeom2018}; Table~\ref{tab:ablations}, the ablations, which include the class-balanced~\cite{cui2019classbalanced}, focal~\cite{lin2017focal}, and label-distribution-aware margin (LDAM)~\cite{cao2019ldam} losses; Table~\ref{tab:backbones}, further backbones, among them the FT-Transformer~\cite{gorishniy2021}; Table~\ref{tab:imaging}, the imaging data, with bias-corrected and accelerated (BCa) bootstrap intervals~\cite{efron1987bca}; Table~\ref{tab:scaling}, the scaling with the number of clients; Table~\ref{tab:replay}, the replay of the logged uploads; and Tables~\ref{tab:dp20} and~\ref{tab:dp100}, the schedules at matched privacy budgets with 20 and 100 clients. The fixed schedule of Table~\ref{tab:replay} is the uniform period $(4,4)$ without warm-up. The fixed periods $K=4$ and $K=8$ of Tables~\ref{tab:dp20} and~\ref{tab:dp100} have a warm-up of five rounds and are then due when $r\bmod K=0$, which gives 19 and 12 releases instead of the 18 and 11 of rule (16) of the article; \tierfed{}-DP follows rule (16). Unless a caption states otherwise, values are means over seeds, $\pm$ denotes the standard deviation over seeds, and, where an entry is in bold, it is the best AUROC of its group.

\begin{table}[t]
\caption{Sweep of the fixed period $(K_{\mathrm s},K_{\mathrm m})$ with local-only rounds (top) and comparison of the exact and the greedy solvers of \tierfed{} at the quotas for which each was run (bottom): final-round AUROC and upload over 60 rounds (10 seeds).}
\label{tab:periods}
\centering\small
\begin{tabular}{@{}lcr@{}}
\toprule
$(K_{\mathrm s},K_{\mathrm m})$ or quota, solver & AUROC & Upload (MB)\\
\midrule
(2,1) & 0.6317 $\pm$ 0.0097 & 170.0 \\
(2,2) & \textbf{0.6362} $\pm$ 0.0110 & 112.5 \\
\midrule
(4,1) & 0.6314 $\pm$ 0.0102 & 145.3 \\
(4,2) & 0.6370 $\pm$ 0.0112 & 87.8 \\
(4,4) & \textbf{0.6407} $\pm$ 0.0114 & 58.1 \\
\midrule
(8,1) & 0.6130 $\pm$ 0.0079 & 132.1 \\
(8,2) & 0.6155 $\pm$ 0.0077 & 74.6 \\
(8,4) & \textbf{0.6254} $\pm$ 0.0073 & 44.9 \\
(8,8) & 0.6250 $\pm$ 0.0088 & 29.0 \\
\midrule
$\rho=0.15$, exact & \textbf{0.6070} $\pm$ 0.0119 & 45.7 \\
\midrule
$\rho=0.25$, exact & \textbf{0.6269} $\pm$ 0.0082 & 65.3 \\
$\rho=0.25$, greedy & \textbf{0.6269} $\pm$ 0.0082 & 65.3 \\
\midrule
$\rho=0.40$, exact & \textbf{0.6366} $\pm$ 0.0119 & 84.9 \\
$\rho=0.40$, greedy & 0.6283 $\pm$ 0.0112 & 95.7 \\
\midrule
$\rho=0.60$, exact & \textbf{0.6290} $\pm$ 0.0093 & 116.2 \\
\midrule
$\rho=0.80$, exact & \textbf{0.6296} $\pm$ 0.0094 & 167.2 \\
\bottomrule
\end{tabular}
\end{table}

\begin{table*}[t]
\caption{Sweeps of the baselines (10 seeds): initial period $\tau_0$ of AdaComm, threshold $\xi$ of LAG, and the methods with an SGD local solver (learning rate $5\times10^{-2}$), among them SCAFFOLD.}
\label{tab:baselines-legacy}
\centering\small
\fit{\begin{tabular}{@{}lccrr@{}}
\toprule
Configuration & AUROC & AUPRC & Upload (MB) & Transmitting rounds\\
\midrule
AdaComm, $\tau_0=2$ & 0.6301 $\pm$ 0.0095 & 0.177 & 109.2 & 30 \\
AdaComm, $\tau_0=4$ & 0.6275 $\pm$ 0.0068 & 0.174 & 54.6 & 15 \\
AdaComm, $\tau_0=8$ & \textbf{0.6302} $\pm$ 0.0071 & 0.168 & 29.1 & 8 \\
AdaComm, $\tau_0=16$ & 0.6125 $\pm$ 0.0111 & 0.152 & 18.2 & 5 \\
\midrule
FedAvg, SGD solver & \textbf{0.6396} $\pm$ 0.0109 & 0.172 & 218.4 & 60 \\
\midrule
FedPer, SGD solver & \textbf{0.6379} $\pm$ 0.0110 & 0.170 & 217.8 & 60 \\
\midrule
LAG, $\xi=10$ & 0.6351 $\pm$ 0.0124 & 0.179 & 218.4 & 60 \\
LAG, $\xi=100$ & 0.6351 $\pm$ 0.0124 & 0.179 & 218.4 & 60 \\
LAG, $\xi=1000$ & \textbf{0.6352} $\pm$ 0.0126 & 0.179 & 168.4 & 60 \\
LAG, $\xi=10000$ & 0.5673 $\pm$ 0.0498 & 0.135 & 6.8 & 7 \\
\midrule
SCAFFOLD, SGD solver & \textbf{0.6383} $\pm$ 0.0106 & 0.173 & 436.9 & 60 \\
\midrule
TierFed, SGD solver & \textbf{0.6251} $\pm$ 0.0095 & 0.157 & 65.3 & 18 \\
\bottomrule
\end{tabular}}
\end{table*}

\begin{table}[t]
\caption{Sensitivity of the privacy planner to its noise weight $\lambda$ (10 seeds): selected interval $K$, which is common to the two shared tiers, calibrated noise multiplier $z$, number of releases, final-round AUROC, and upload. Top: 20 clients at $q=1$, $\varepsilon^\star=4$. Bottom: 100 clients at $q=0.2$, $\varepsilon^\star=4$.}
\label{tab:lambda}
\centering\small
\begin{tabular}{@{}lrrrcr@{}}
\toprule
$\lambda$ & $K$ & $z$ & Releases & AUROC & Upload (MB)\\
\midrule
0.3 & 2 & 7.33 & 32 & 0.5237 $\pm$ 0.0150 & 116.2 \\
1.0 & 3 & 6.22 & 23 & 0.5211 $\pm$ 0.0121 & 83.5 \\
3.0 & 4 & 5.50 & 18 & \textbf{0.5318} $\pm$ 0.0081 & 65.3 \\
10.0 & 5 & 5.18 & 16 & 0.5200 $\pm$ 0.0071 & 58.1 \\
\midrule
0.3 & 2 & 1.85 & 32 & \textbf{0.5158} $\pm$ 0.0256 & 116.3 \\
1.0 & 3 & 1.67 & 23 & 0.5122 $\pm$ 0.0226 & 83.0 \\
3.0 & 4 & 1.56 & 18 & 0.5084 $\pm$ 0.0192 & 64.5 \\
10.0 & 6 & 1.47 & 14 & 0.5114 $\pm$ 0.0283 & 50.8 \\
\bottomrule
\end{tabular}
\end{table}

\begin{table}[t]
\caption{Exposure of the transmitted updates without DP (10 seeds): parameters exposed per transmitting round, normalized mean squared error (NMSE) and feature correlation of an update-inversion attack on the slices that were actually transmitted, AUC of a membership-inference attack (MIA) on the deployed models, for which 0.5 corresponds to no leakage, and upload. Arrows give the direction that is better for privacy.}
\label{tab:exposure}
\centering\small
\fit{\begin{tabular}{@{}lrrccr@{}}
\toprule
Method & \shortstack{Params\\exposed} & NMSE $\uparrow$ & Corr.\ $\downarrow$ & \shortstack{MIA\\AUC $\downarrow$} & \shortstack{Upload\\(MB)}\\
\midrule
AdaComm & 45,506 & 11.58 & 0.065 & 0.534 & 29.1 \\
FedAvg & 45,506 & 9.04 & 0.103 & 0.534 & 218.4 \\
FedPer & 45,376 & 9.04 & 0.088 & 0.535 & 217.8 \\
Fixed schedule & 27,509 & 11.93 & 0.048 & 0.526 & 82.9 \\
TierFed & 45,376 & 10.45 & 0.089 & 0.540 & 65.3 \\
\bottomrule
\end{tabular}}
\end{table}

\begin{table}[t]
\caption{Ablations (10 seeds). Top: private head and schedule, and \tierfed{} with and without the private head. Middle: length $T_w$ of the warm-up. Bottom: loss used against class imbalance, for FedPer and \tierfed{} (CE: cross-entropy; LDAM: label-distribution-aware margin).}
\label{tab:ablations}
\centering\small
\fit{\begin{tabular}{@{}lccr@{}}
\toprule
Configuration & AUROC & AUPRC & Upload (MB)\\
\midrule
FedAvg: shared head, every round & 0.6203 $\pm$ 0.0100 & 0.171 & 218.4 \\
FedPer: private head, every round & 0.6215 $\pm$ 0.0087 & 0.171 & 217.8 \\
Fixed schedule, private head & \textbf{0.6339} $\pm$ 0.0100 & 0.177 & 82.9 \\
Fixed schedule, shared head & 0.6336 $\pm$ 0.0102 & 0.177 & 83.5 \\
TierFed, private head & 0.6269 $\pm$ 0.0082 & 0.175 & 65.3 \\
TierFed, shared head & 0.6264 $\pm$ 0.0073 & 0.174 & 65.5 \\
\midrule
TierFed, $T_w=0$ & 0.6270 $\pm$ 0.0080 & 0.173 & 54.5 \\
TierFed, $T_w=2$ & 0.6323 $\pm$ 0.0092 & 0.176 & 58.1 \\
TierFed, $T_w=5$ & 0.6269 $\pm$ 0.0082 & 0.175 & 65.3 \\
TierFed, $T_w=10$ & 0.6289 $\pm$ 0.0093 & 0.176 & 79.9 \\
TierFed, $T_w=20$ & \textbf{0.6354} $\pm$ 0.0112 & 0.179 & 108.9 \\
\midrule
Class-balanced, FedPer & 0.6232 $\pm$ 0.0108 & 0.173 & 217.8 \\
Class-balanced, TierFed & \textbf{0.6264} $\pm$ 0.0071 & 0.173 & 65.3 \\
\midrule
Focal, FedPer & 0.6198 $\pm$ 0.0100 & 0.170 & 217.8 \\
Focal, TierFed & \textbf{0.6237} $\pm$ 0.0089 & 0.173 & 65.3 \\
\midrule
LDAM, FedPer & 0.6150 $\pm$ 0.0114 & 0.164 & 217.8 \\
LDAM, TierFed & \textbf{0.6276} $\pm$ 0.0091 & 0.171 & 65.3 \\
\midrule
Inverse-frequency CE, FedPer & 0.6215 $\pm$ 0.0087 & 0.171 & 217.8 \\
Inverse-frequency CE, TierFed & \textbf{0.6269} $\pm$ 0.0082 & 0.175 & 65.3 \\
\bottomrule
\end{tabular}}
\end{table}

\begin{table}[t]
\caption{Backbones on the readmission data (5 seeds): final-round AUROC, AUPRC, upload, and number of parameters; the transformer runs used a fused attention kernel that is not deterministic.}
\label{tab:backbones}
\centering\small
\setlength{\tabcolsep}{4pt}\fit{\begin{tabular}{@{}lccrr@{}}
\toprule
Backbone, method & AUROC & AUPRC & \shortstack{Upload\\(MB)} & Params\\
\midrule
MLP / FedAvg & 0.6236 $\pm$ 0.0103 & 0.171 & 218.4 & 45,506 \\
MLP / FedPer & 0.6240 $\pm$ 0.0087 & 0.172 & 217.8 & 45,506 \\
MLP / Fixed schedule & \textbf{0.6373} $\pm$ 0.0117 & 0.176 & 82.9 & 45,506 \\
MLP / TierFed & 0.6298 $\pm$ 0.0128 & 0.175 & 65.3 & 45,506 \\
\midrule
ResNet / FedAvg & 0.6258 $\pm$ 0.0192 & 0.168 & 2007.3 & 418,178 \\
ResNet / FedPer & 0.6137 $\pm$ 0.0107 & 0.161 & 2006.0 & 418,178 \\
ResNet / Fixed schedule & \textbf{0.6382} $\pm$ 0.0150 & 0.177 & 1006.8 & 418,178 \\
ResNet / TierFed & 0.6279 $\pm$ 0.0130 & 0.175 & 603.5 & 418,178 \\
\midrule
FT-Transformer / FedAvg & \textbf{0.6401} $\pm$ 0.0134 & 0.177 & 581.8 & 121,218 \\
FT-Transformer / FedPer & 0.6352 $\pm$ 0.0115 & 0.173 & 581.2 & 121,218 \\
FT-Transformer / Fixed schedule & 0.6386 $\pm$ 0.0116 & 0.172 & 270.7 & 121,218 \\
FT-Transformer / TierFed & 0.6345 $\pm$ 0.0083 & 0.171 & 165.2 & 121,218 \\
\bottomrule
\end{tabular}}
\end{table}

\begin{table*}[t]
\caption{Imaging data (PneumoniaMNIST, Dirichlet label skew over 20 clients, of which 16 to 20 receive data, 10 seeds): final-round AUROC with its bias-corrected and accelerated (BCa) bootstrap 95\% interval, AUPRC, upload, transmitting rounds (Tx), and $p$-value of a one-sided paired $t$-test of noninferiority to FedPer on the AUROC (NI $p$; margin 0.01; this test is not the signed-rank test of the article). Between 3 and 9 clients per seed have a test split with a single class and are excluded identically for every method.}
\label{tab:imaging}
\centering\small
\fit{\begin{tabular}{@{}lccrrc@{}}
\toprule
Method & AUROC [95\% CI] & AUPRC & Upload (MB) & Tx & NI $p$\\
\midrule
FedAvg & 0.9959 [0.9944, 0.9972] & 0.991 & 1214.8 & 60 & 0.037 \\
FedPer & 0.9643 [0.8952, 0.9917] & 0.986 & 1213.7 & 60 & (ref.) \\
FedLAMA-style & \textbf{0.9965} [0.9953, 0.9975] & 0.995 & 877.4 & 60 & 0.036 \\
Fixed schedule & 0.9642 [0.8938, 0.9934] & 0.989 & 614.8 & 37 & \ensuremath{<10^{-4}} \\
TierFed & 0.9698 [0.9536, 0.9820] & 0.984 & 364.1 & 18 & 0.182 \\
\bottomrule
\end{tabular}}
\end{table*}

\begin{table*}[t]
\caption{Scaling with the number of clients (readmission data, Dirichlet $\alpha=1$, 30 rounds, 3 seeds): AUROC, total upload of all clients, wall-clock seconds per round (time of the whole run, evaluation included, divided by the 30 rounds), and peak graphics processing unit (GPU) memory allocated in megabytes.}
\label{tab:scaling}
\centering\small
\fit{\begin{tabular}{@{}lcrrr@{}}
\toprule
Clients, method & AUROC & Upload (MB) & s/round & Peak GPU MB\\
\midrule
10 clients, FedAvg & \textbf{0.6127} $\pm$ 0.0206 & 54.6 & 3.07 & 80 \\
10 clients, FedLAMA-style & 0.6056 $\pm$ 0.0117 & 43.1 & 3.03 & 80 \\
10 clients, FedPer & 0.5935 $\pm$ 0.0072 & 54.5 & 3.03 & 80 \\
10 clients, TierFed & 0.6118 $\pm$ 0.0100 & 20.0 & 3.01 & 80 \\
\midrule
20 clients, FedAvg & 0.6187 $\pm$ 0.0110 & 109.2 & 3.12 & 87 \\
20 clients, FedLAMA-style & 0.6187 $\pm$ 0.0183 & 85.4 & 3.09 & 87 \\
20 clients, FedPer & 0.6189 $\pm$ 0.0235 & 108.9 & 3.11 & 87 \\
20 clients, TierFed & \textbf{0.6205} $\pm$ 0.0242 & 39.9 & 3.04 & 87 \\
\midrule
50 clients, FedAvg & 0.6128 $\pm$ 0.0224 & 273.0 & 3.40 & 108 \\
50 clients, FedLAMA-style & 0.6108 $\pm$ 0.0203 & 219.3 & 3.33 & 108 \\
50 clients, FedPer & \textbf{0.6188} $\pm$ 0.0240 & 272.3 & 3.37 & 108 \\
50 clients, TierFed & 0.6048 $\pm$ 0.0170 & 99.8 & 3.21 & 108 \\
\midrule
100 clients, FedAvg & \textbf{0.6068} $\pm$ 0.0210 & 546.1 & 3.99 & 142 \\
100 clients, FedLAMA-style & 0.6061 $\pm$ 0.0230 & 437.1 & 3.86 & 142 \\
100 clients, FedPer & 0.5929 $\pm$ 0.0255 & 544.5 & 3.86 & 142 \\
100 clients, TierFed & 0.5883 $\pm$ 0.0157 & 199.7 & 3.57 & 142 \\
\bottomrule
\end{tabular}}
\end{table*}

\begin{table*}[t]
\caption{Replay of the logged uploads over TCP for the runs of a wide MLP (1,740,290 parameters, 3 seeds; cell-edge profile for one seed). The per-client uploads of each run are sent over loopback TCP paced at the rate of the link profile. The table gives the total replayed upload time, the prediction of the model of Section~\mbox{III-B}, their relative difference (prediction minus replay, as a percentage of the replay), the total compute time of the slowest client, and the time to an AUROC of 0.60 and of 0.62, computed as the modeled upload time plus the measured compute time (``never'': not reached in 60 rounds). The profile ``5G NR, typical'' (50~Mb/s, 20~ms, 0.2~W) is not used in the article.}
\label{tab:replay}
\centering\small
\fit{\begin{tabular}{@{}llrrrrrr@{}}
\toprule
Profile & Method & Replayed (s) & Model (s) & Diff.\ (\%) & Compute (s) & $t_{0.60}$ (s) & $t_{0.62}$ (s)\\
\midrule
Wired WAN & AdaComm & 5.2 & 5.1 & \ensuremath{-}1.7 & 10 & 1 & never \\
Wired WAN & FedAvg & 34.6 & 34.0 & \ensuremath{-}1.7 & 10 & 1 & 2 \\
Wired WAN & FedPer & 34.6 & 34.0 & \ensuremath{-}1.6 & 10 & 1 & 2 \\
Wired WAN & Fixed schedule & 9.2 & 9.1 & \ensuremath{-}1.7 & 10 & 1 & 2 \\
Wired WAN & TierFed & 10.4 & 10.2 & \ensuremath{-}1.9 & 10 & 1 & 2 \\
\midrule
5G NR, typical & AdaComm & 10.2 & 10.2 & \ensuremath{-}0.2 & 10 & 1 & never \\
5G NR, typical & FedAvg & 68.2 & 68.0 & \ensuremath{-}0.2 & 10 & 2 & 3 \\
5G NR, typical & FedPer & 68.1 & 68.0 & \ensuremath{-}0.2 & 10 & 3 & 4 \\
5G NR, typical & Fixed schedule & 18.2 & 18.1 & \ensuremath{-}0.2 & 10 & 1 & 4 \\
5G NR, typical & TierFed & 20.4 & 20.4 & \ensuremath{-}0.2 & 10 & 3 & 4 \\
\midrule
5G NR, cell edge & AdaComm & 100.5 & 100.5 & \ensuremath{+}0.1 & 10 & 11 & never \\
5G NR, cell edge & FedAvg & 669.6 & 670.1 & \ensuremath{+}0.1 & 11 & 13 & 24 \\
5G NR, cell edge & FedPer & 669.2 & 669.7 & \ensuremath{+}0.1 & 10 & 23 & 34 \\
5G NR, cell edge & Fixed schedule & 178.5 & 178.6 & \ensuremath{+}0.1 & 10 & 12 & 35 \\
5G NR, cell edge & TierFed & 200.8 & 200.9 & \ensuremath{+}0.1 & 10 & 23 & 34 \\
\bottomrule
\end{tabular}}
\end{table*}

\clearpage
% \clearpage with pending table* floats leaves the column-room register of the float page; reset it for this page.
\makeatletter\@floatplacement\global\@colroom\@colht\global\vsize\@colht\makeatother
\begin{table}[!t]
\caption{Schedules at matched privacy budgets on the 20 specialty groups for all sampling rates, $q=1$ included (10 seeds, $C=1$, $\delta=10^{-5}$): calibrated noise multiplier $z$, number of releases, and final-round AUROC for each sampling rate $q$ and budget $\varepsilon^\star$.}
\label{tab:dp20}
\centering\small\renewcommand{\arraystretch}{0.90}
\fit{\begin{tabular}{@{}lrrc@{}}
\toprule
$q$, $\varepsilon^\star$, method & $z$ & Releases & AUROC\\
\midrule
0.5, 1, FedAvg & 19.14 & 60 & 0.5073 $\pm$ 0.0163 \\
0.5, 1, FedPer & 19.14 & 60 & 0.5102 $\pm$ 0.0092 \\
0.5, 1, Fixed $K=4$ & 10.95 & 19 & 0.5150 $\pm$ 0.0083 \\
0.5, 1, Fixed $K=8$ & 8.83 & 12 & \textbf{0.5208} $\pm$ 0.0086 \\
0.5, 1, TierFed-DP & 12.00 & 23 & 0.5115 $\pm$ 0.0094 \\
\midrule
0.5, 2, FedAvg & 9.84 & 60 & 0.5079 $\pm$ 0.0193 \\
0.5, 2, FedPer & 9.84 & 60 & 0.5090 $\pm$ 0.0092 \\
0.5, 2, Fixed $K=4$ & 5.72 & 19 & 0.5178 $\pm$ 0.0149 \\
0.5, 2, Fixed $K=8$ & 4.66 & 12 & \textbf{0.5248} $\pm$ 0.0113 \\
0.5, 2, TierFed-DP & 6.25 & 23 & 0.5118 $\pm$ 0.0105 \\
\midrule
0.5, 4, FedAvg & 5.18 & 60 & 0.5104 $\pm$ 0.0198 \\
0.5, 4, FedPer & 5.18 & 60 & 0.5135 $\pm$ 0.0089 \\
0.5, 4, Fixed $K=4$ & 3.09 & 19 & 0.5218 $\pm$ 0.0134 \\
0.5, 4, Fixed $K=8$ & 2.57 & 12 & \textbf{0.5292} $\pm$ 0.0110 \\
0.5, 4, TierFed-DP & 3.36 & 23 & 0.5127 $\pm$ 0.0106 \\
\midrule
0.5, 8, FedAvg & 2.84 & 60 & 0.5148 $\pm$ 0.0173 \\
0.5, 8, FedPer & 2.84 & 60 & 0.5162 $\pm$ 0.0107 \\
0.5, 8, Fixed $K=4$ & 1.76 & 19 & 0.5300 $\pm$ 0.0163 \\
0.5, 8, Fixed $K=8$ & 1.48 & 12 & \textbf{0.5460} $\pm$ 0.0164 \\
0.5, 8, TierFed-DP & 1.89 & 23 & 0.5144 $\pm$ 0.0083 \\
\midrule
1.0, 4, FedAvg & 10.04 & 60 & 0.5100 $\pm$ 0.0192 \\
1.0, 4, FedPer & 10.04 & 60 & 0.5174 $\pm$ 0.0103 \\
1.0, 4, Fixed $K=4$ & 5.65 & 19 & \textbf{0.5294} $\pm$ 0.0083 \\
1.0, 4, Fixed $K=8$ & 4.49 & 12 & 0.5272 $\pm$ 0.0104 \\
1.0, 4, TierFed-DP & 6.22 & 23 & 0.5211 $\pm$ 0.0121 \\
\bottomrule
\end{tabular}}
\end{table}

\begin{table}[!t]
\caption{Schedules at matched privacy budgets on a federation of 100 clients with Dirichlet label skew (10 seeds, $C=1$, $\delta=10^{-5}$): calibrated noise multiplier $z$, number of releases, and final-round AUROC for each sampling rate $q$ and budget $\varepsilon^\star$.}
\label{tab:dp100}
\centering\small\renewcommand{\arraystretch}{0.90}
\fit{\begin{tabular}{@{}lrrc@{}}
\toprule
$q$, $\varepsilon^\star$, method & $z$ & Releases & AUROC\\
\midrule
0.2, 1, FedAvg & 7.86 & 60 & 0.5062 $\pm$ 0.0197 \\
0.2, 1, FedPer & 7.86 & 60 & 0.5069 $\pm$ 0.0115 \\
0.2, 1, Fixed $K=4$ & 4.73 & 19 & 0.5109 $\pm$ 0.0203 \\
0.2, 1, Fixed $K=8$ & 3.95 & 12 & 0.5112 $\pm$ 0.0148 \\
0.2, 1, TierFed-DP & 5.12 & 23 & \textbf{0.5114} $\pm$ 0.0219 \\
\midrule
0.2, 2, FedAvg & 4.15 & 60 & 0.5126 $\pm$ 0.0216 \\
0.2, 2, FedPer & 4.15 & 60 & 0.5076 $\pm$ 0.0117 \\
0.2, 2, Fixed $K=4$ & 2.64 & 19 & 0.5108 $\pm$ 0.0206 \\
0.2, 2, Fixed $K=8$ & 2.28 & 12 & \textbf{0.5139} $\pm$ 0.0188 \\
0.2, 2, TierFed-DP & 2.83 & 23 & 0.5100 $\pm$ 0.0246 \\
\midrule
0.2, 4, FedAvg & 2.31 & 60 & 0.5159 $\pm$ 0.0200 \\
0.2, 4, FedPer & 2.31 & 60 & 0.5064 $\pm$ 0.0147 \\
0.2, 4, Fixed $K=4$ & 1.59 & 19 & 0.5117 $\pm$ 0.0201 \\
0.2, 4, Fixed $K=8$ & 1.42 & 12 & \textbf{0.5180} $\pm$ 0.0232 \\
0.2, 4, TierFed-DP & 1.67 & 23 & 0.5122 $\pm$ 0.0226 \\
\midrule
0.2, 8, FedAvg & 1.39 & 60 & 0.5175 $\pm$ 0.0211 \\
0.2, 8, FedPer & 1.39 & 60 & 0.5064 $\pm$ 0.0142 \\
0.2, 8, Fixed $K=4$ & 1.04 & 19 & 0.5112 $\pm$ 0.0204 \\
0.2, 8, Fixed $K=8$ & 0.93 & 12 & \textbf{0.5217} $\pm$ 0.0216 \\
0.2, 8, TierFed-DP & 1.09 & 23 & 0.5101 $\pm$ 0.0220 \\
\midrule
0.5, 4, FedAvg & 5.18 & 60 & 0.5196 $\pm$ 0.0203 \\
0.5, 4, FedPer & 5.18 & 60 & 0.5135 $\pm$ 0.0094 \\
0.5, 4, Fixed $K=4$ & 3.09 & 19 & 0.5304 $\pm$ 0.0151 \\
0.5, 4, Fixed $K=8$ & 2.57 & 12 & \textbf{0.5421} $\pm$ 0.0159 \\
0.5, 4, TierFed-DP & 3.36 & 23 & 0.5274 $\pm$ 0.0281 \\
\bottomrule
\end{tabular}}
\end{table}

\section{Released Data}
\label{supp:data}
The final metrics of every run are in \texttt{results/\allowbreak tables/\allowbreak collected\_all\_raw.csv} of the released repository, and the per-round histories and per-round bit counts of every federated run are in \texttt{results/logs}; the 20 reference-bound runs have a summary only. The repository held 2,760 run logs when this document was built: the 946 runs of Table~\ref{tab:families}, the single-seed check of Section~\ref{supp:sync}, and 1,813 runs of the experiments of Section~\ref{supp:legacy} and of earlier experiments that are not reported.

\bibliographystyle{IEEEtran}
\bibliography{references}
\end{document}
